% Cluster A: comparison with standard continuum-hull constructions
% Review-required new expository material; promotes no manuscript claim.
% Body-only LaTeX; preamble macros, cite keys, and cross-cluster labels per CONVENTIONS.md.

\section{Comparison with standard continuum-hull constructions}
\label{sec:comparison}

The certificate formalism developed from Section~\ref{sec:certificates}
onward is offered as a replacement for a family of standard constructions,
and it is fair to ask what, exactly, is being replaced and why.  This
section compares the framework with five natural comparison points:
least or minimal connected full hulls; monotone
absorption of pulled excess; regulated and Hubbard trees in uniformized or
locally connected coordinates; Moore decompositions and monotone
quotients; and the lamination-lift candidate of the source manuscript.
The two obstructions of Section~\ref{sec:obstructions} already dispose of
the first two; the remaining three fail for structural reasons that are
worth stating precisely, because they explain several design decisions of
the certificate formalism that would otherwise look pedantic.  Throughout,
the comparison is with respect to the task fixed in the introduction:
join finitely many pairwise disjoint full atoms inside a Jordan
localization, avoid a frozen forbidden set, and \emph{preserve designated
prime-end accesses together with their cyclic order}.  Constructions that
solve a different problem well are not thereby criticized; the claim is
only that none of the constructions, in the forms reviewed here, records
all of these requirements simultaneously, and that the last of them ---
the lamination lift --- is an open candidate, not a construction.

\subsection{Least and minimal connected full hulls}
\label{subsec:A-comp-least}

The most tempting rule is extremal: ``take the least connected full hull
of the required package.''  Theorem~\ref{thm:no-least} (manuscript
Theorem~6.12) shows the rule names no object, already for the two-point
package $S=\{-1,1\}$ with identity dynamics and no obstacles: the upper
and lower semicircles are carriers meeting exactly in $S$, so a least
carrier would be contained in their disconnected intersection.  The
failure is specific to \emph{leastness}: minimal carriers exist in
abundance --- every arc from $-1$ to $1$ is one --- and form a large
antichain.  An extremal selection principle therefore cannot even break
the tie between the two semicircles, and the tie is not innocent: the two
choices pass on opposite sides of the origin, and in any configuration
with obstacles or relative protected collar tails that choice of side is visible to
every later routing step.  The source manuscript sharpens the point
dynamically \cite[Proposition~6.11]{MWman}: a package, full connected
carriers, critical-value-free degree-one selected-carrier dynamics, and
identical coarse orbit labels are realized simultaneously by an interval
and by a triod, and attachment-freeness holds in only one of the two
models, so no data of the listed kinds can select between them.

The connected-hull descriptions in the renormalization
literature are not at fault here, and the comparison should not be read
as a criticism of them.  Dudko's roadmap describes virtual Julia sets as
connected hulls of Cantor families of small Julia sets carrying a finite
part of the critical orbit \cite{Dud25}; that description specifies
\emph{which} continua a carrier must contain, and it does not purport to
specify \emph{how} they are joined.  Theorem~\ref{thm:no-least} shows
that the routing is genuinely additional data --- in the manuscript's
words, a regulated connector system, together with its access or isotopy
class and a forward-coherence proof, is additional data \cite{MWman} ---
and the certificate of Definition~\ref{def:certificate} is exactly a
format for that data, with Theorem~\ref{thm:free-component} deciding when
the format can be filled in at all.

\subsection{Monotone absorption of pulled excess}
\label{subsec:A-comp-absorption}

The second rule is completion: accept any proposed carrier and absorb,
monotonically, whatever excess the dynamics pulls back onto it, until the
process stabilizes.  Unlike the least hull, this rule is well defined:
the manuscript proves that the finite inflationary completion operator
stabilizes at a least fixed point, unique and conjugation-equivariant
\cite[Theorem~6.14]{MWman}.  The defect is not existence but what
completion destroys.  Proposition~\ref{prop:access-collision} (manuscript
Proposition~6.15) exhibits a conformal, critical-value-free, degree-one
branch of $z\mapsto z^2$ in which the very first absorption of pulled
excess --- the relative upper half-disk attaching to the interval along
its entire length --- erases a protected prime-end access and its
conformally transported copy.  Monotonicity, the rule's defining virtue,
is precisely what makes the loss permanent: enlargement can only add
points, and no addition restores a complement access through a set
already absorbed.  Moreover, every coordinate of the manuscript's
collision bookkeeping except the access coordinate vanishes in this
example, so no amount of critical-value, degree, or sheet bookkeeping
detects the erasure (Section~\ref{subsec:B-moral}).

The certificate formalism keeps monotone completion but reverses the
order of commitments.  The relative protected collar tails are frozen into the
inventory $Q$ \emph{before} any absorption is contemplated; a proposed
absorption that would enter such a tail is not a step of a completion
but a named violation, and the persistence clause of
Theorem~\ref{thm:free-component}, together with the audit of
Theorem~\ref{thm:audit}, converts each such violation into a certificate
--- a merger, a port loss, a separation --- that survives every further
monotone enlargement.  Absorption remains available for everything
unprotected; what changes is that ``protected'' is a constraint checked
by the plane, not an intention held by the author.

\subsection[Regulated and Hubbard trees]
{Regulated trees in uniformized coordinates and Hubbard trees}
\label{subsec:A-comp-trees}

For a postcritically finite polynomial $P$ with connected filled Julia
set, the classical Hubbard tree is the regulated hull in $K(P)$ of the
finite union of the forward critical orbits \cite{DHOrsay}.  Arbitrary finite
invariant sets, or local connectivity alone, do not define the classical
Hubbard tree.  The construction uses regulated arcs, uniformized
coordinates, and canonical internal arcs.  Local connectivity gives a
continuous extension of the exterior Riemann map to the boundary; that
extension is a homeomorphism only when the boundary is a Jordan curve
\cite{DHOrsay,Pom92}.  Kahn--Lyubich impose finite-decoration combinatorics
and prove a~priori bounds for that class, using puzzle pieces, rays, and
equipotentials in the construction \cite{KL08}.

None of these steps is guaranteed by the standing hypotheses of the
present paper.  The atoms are arbitrary full continua; no Riemann map is
asserted to extend, no external or internal ray is asserted to land, and
regulated arcs need not exist.  This is not a hypothetical worry but the
operating regime of the intended applications, in which local
connectivity of the relevant continua is exactly what is not known.  The
closest publicly available construction is instructive as a boundary
marker.  In its stated primitive-renormalization setting, the preprint of
Kahn--Kapiamba--Lyubich uses a rational-lamination quotient, a locally
connected combinatorial model, and the associated natural projection
\cite[\S\S2.5--2.6]{KKL26}.  This is not a general construction for
arbitrary full continua.  As the source manuscript's own boundary records
\cite{MWman}, the preprint is a positive control in its
primitive-renormalization category, not a theorem for the arbitrary
packages contemplated here.  The framework of this paper is
calibrated to the complementary regime: Theorem~\ref{thm:realization}(6)
uses no local connectivity of any atom, the tree of the certificate is
produced by the exact criterion of Theorem~\ref{thm:free-component} from
planar data alone, and no canonicity of the spanning tree is claimed
anywhere --- the tree is data, chosen in a free component whose existence
is the actual content.  What is lost relative to the Hubbard-tree regime
is canonicity; what is gained is existence without a locally connected
model whenever the exact free-component criterion holds.

\subsection{Moore decompositions and monotone quotients}
\label{subsec:A-comp-moore}

A third classical route collapses instead of routing.  Moore's theorem
states that if $\mathcal G$ is an upper-semicontinuous decomposition of
$S^2$ into nonseparating continua and the quotient has more than one
point, then $S^2/\mathcal G$ is homeomorphic to $S^2$ \cite{Moo25}.
One is therefore tempted to collapse every atom (or every
completed packet) to a point, join the finitely many image points by
arcs in the quotient sphere --- where the joining problem is trivial ---
and pull the picture back.  For the connection problem \emph{without}
protected data this can be a serviceable device and is one standard
mechanism for passing to a simpler quotient model.

For the access-preserving problem the unmarked quotient is not, by itself,
a certificate.  Collapsing an atom sends all of its points to one quotient
point, and neither that point nor a tree incident to it records which named
prime ends, selected sides, or relative collar tails are to be protected.
Crosscuts in the complement can have quotient images, and in favorable
decompositions distinct approaches may still be recoverable; the point is
that Moore's theorem alone supplies none of the required marking or
survival conclusions.  Those data would have to be transported separately.

Likewise, the preimage of an arc under a continuous monotone quotient is
a continuum, but generally a saturated one.  It does not automatically
have singleton attachment traces, retain selected prime-end accesses, or
preserve marking data; a fiber can contain more attachment material or
meet a forbidden collar.  Thus a Moore quotient may support a construction only
after it is augmented by the same access and incidence data that the
certificate formalism records explicitly.  The present realization works
in the original plane, carries protected accesses as crosscut chains with
selected sides, and attaches its connector at declared free ports with
singleton traces (Theorem~\ref{thm:realization}(2)).

\subsection{The lamination-lift candidate}
\label{subsec:A-comp-lamination}

The most serious competitor is the manuscript's own candidate, recorded
in its Section~10.3 \cite{MWman}, and the comparison must be exact:
the candidate is presented there, and here, as a candidate only.
Suppose, following the non-locally-connected construction of
Kahn--Kapiamba--Lyubich \cite{KKL26}, that an ambient carrier admits a
monotone semiconjugacy
\[
\pi\colon K\longrightarrow K_{\mathrm{com}}
\]
onto a locally connected combinatorial model, together with a
forward-coherent regulated tree $T_{\mathrm{com}}\subset K_{\mathrm{com}}$
containing the prescribed finite marked set.  The canonical candidate connector is
\begin{equation}\label{eq:A-lam}
H_{\mathrm{lam}}\;=\;\Fill\bigl(\pi^{-1}(T_{\mathrm{com}})\bigr),
\end{equation}
equation~(9.3) of \cite{MWman}.  Part of what is wanted is genuinely
true, by soft arguments: a monotone preimage of a compact connected set
is compact and connected, so \eqref{eq:A-lam} is a full continuum
containing the preimage of the prescribed marked set; the semiconjugacy provides forward
coherence before fullification; and fullification preserves genuine
polynomial forward invariance.

The missing assertion is not connectedness.  Measured against the
certificate clauses, two failure channels are named in the manuscript,
and both remain unresolved.  First, a projection fiber
$\pi^{-1}(x)$ over a point of the tree, or a component added by
fullification, may enter a protected collar otherwise than at a declared
free singleton port; this is the erasure mechanism of
Proposition~\ref{prop:access-collision} wearing different clothes, and
it is instructive that the manuscript rejects the upper-half-disk
collision at exactly this clause, the excess occupying the protected
upper collar with an interval attachment trace \cite{MWman}.  Second, a
disconnected stage package may still split across iterated pullback
components, so that containing the prescribed marked set does not yet realize
the complete package routing.  The manuscript collects the clauses a
proof would have to establish --- authorized localization and complete
frozen inventory, collar avoidance except at declared free singleton
ports, realization of every package component and route, obstacle
handling inside the localization, complete access marking, and an
incidence bound or compression --- as its Open Problem~10.6 (Virtual
Protected Routing Certificate).  That problem is open, and nothing in
this paper resolves any part of it; consistently with the manuscript's
Warning~10.59, the candidate \eqref{eq:A-lam} is discussed here only to
locate it on the map.

What this paper contributes to the candidate is the deciding instrument,
not a verdict.  For any frozen inventory that a future argument
supplies, Theorem~\ref{thm:free-component} decides certificate existence
exactly; Theorems~\ref{thm:common-domain} and~\ref{thm:audit} and
Proposition~\ref{prop:cellwise} make the decision checkable packet by
packet; and the two failure channels above fall under the typed,
monotone-persistent defect certificates of Section~\ref{sec:audit} ---
a fiber in a collar is a port-loss or merger defect, a split package is
a separation in the sense of the persistence clause of
Theorem~\ref{thm:free-component}.  Whether the candidate passes the test
is precisely what is open.

\subsection{Summary}
\label{subsec:A-comp-summary}

Table~\ref{tab:A-comparison} compresses the comparison.  The pattern is
uniform: each standard construction either fails to exist (the least
hull), can destroy the protected data (absorption), does not by itself
record those data (Moore quotients),
requires the hypothesis whose absence is the whole difficulty (regulated
and Hubbard trees), or is an explicitly open candidate (the lamination
lift).  The certificate framework renounces canonicity --- the connector
is data, not an invariant --- and in exchange obtains, for arbitrary
full atoms and with no local connectivity of the atoms: realization of every
supplied certificate with
access survival and cyclic-order preservation
(Theorem~\ref{thm:realization}); an exact existence criterion on frozen
data (Theorem~\ref{thm:free-component}); checkable sufficient tests
(Theorem~\ref{thm:common-domain}, Proposition~\ref{prop:cellwise});
order-independent completion with a three-clause audit
(Theorem~\ref{thm:audit}); and failure certificates that persist under
monotone enlargement.  The effect is to convert ``choose a hull,''
an instruction that Theorem~\ref{thm:no-least} shows to be ill posed and
Proposition~\ref{prop:access-collision} shows to be destructive, into a
checkable planar test with a stable answer in both directions.  This is
the precise sense in which the contribution of the paper is proof
technology: no new theorem of classical planar topology is claimed, and
every statement about the open constructions touched on above ---
virtual carriers, lamination lifts, and their dynamical certificates ---
remains exactly as open here as in the source manuscript.

\begin{table}[H]
\centering
\small
\begin{tabularx}{\textwidth}{@{}P{0.25\textwidth}P{0.22\textwidth}P{0.22\textwidth}Y@{}}
\toprule
Construction & Existence & Access guarantee & Failure witness\\
\midrule
least connected full hull & no (Thm.~\ref{thm:no-least}) & --- & ---\\
monotone absorption & least fixed point \cite{MWman} &
  erasable (Prop.~\ref{prop:access-collision}) & post hoc only\\
regulated/Hubbard tree & needs locally connected model \cite{DH85,KKL26} &
  native category only & ---\\
Moore monotone quotient & specified nontrivial u.s.c.\ decomposition
  into nonseparating continua \cite{Moo25} &
  markings not automatic & ---\\
lamination lift \eqref{eq:A-lam} & candidate \cite{MWman} & open & open\\
protected certificate & exact criterion (Thm.~\ref{thm:free-component}) &
  preserved (Thm.~\ref{thm:realization}) &
  persistent (Thm.~\ref{thm:audit})\\
\bottomrule
\end{tabularx}
\caption{The five standard constructions against the certificate
framework, measured on the access-preserving connection task.  ``Existence''
asks whether the construction produces a well-defined carrier
for arbitrary pairwise disjoint full atoms in a Jordan localization;
``access guarantee'' asks whether designated prime-end data survive
with their cyclic order; ``failure witness'' asks whether a negative
outcome is witnessed by a certificate that persists under monotone
enlargement.  The lamination-lift row records candidacy, not failure:
its clauses are the manuscript's open problem, and no verdict is implied
in either direction.}
\label{tab:A-comparison}
\end{table}
